\section{Введение}
\label{sec:Introduction}

Визуализация трёхмерных поверхностей — одна из фундаментальных задач компьютерной графики. Во многом именно она определяет комфорт от работы с системами автоматизированного проектирования, медицинскими тренажерами и изображениями, средствами трехмерного моделирования и другими интерактивными приложениями. Стремящиеся к бесконечной детализации современные трёхмерные модели, получающиеся в результате трёхмерных реконструкций и нейронных генераций, становится все сложнее рендерить в реальном времени.

Исторически наиболее распространенным методом представления геометрии является сетка полигонов (mesh), заданная вершинами и способом их объединения в геометрические примитивы. За время своего существования этот метод успел обрости сильной аппаратной поддержкой со стороны графических процессоров. У этого подхода есть ряд критических ограничений: артефакты дискретизации при сильном приближении камеры, экспоненциальный рост затрат видеопамяти и высокая вычислительная сложность при необходимости изменения топологии поверхности (например, при симуляции жидкостей и прочих деформациях).

Альтернативой явному заданию геометрии является использование неявных поверхностей. Для трёхмерных объектов в качестве неявной поверхности зачастую используют нулевую изоповерхнсоть знаковой функции расстояния (Signed Distance Field, SDF). SDF описывает объект как непрерывное поле дистанций, что обеспечивает потенциально бесконечную детализацию и позволяет выполнять сложные геометрические операции (булевы вычитания, сглаживания) с помощью элементарных математических преобразований. Современные исследования в области компактных и нейронных представлений SDF демонстрируют возможность описания сцен при объёмах памяти в десятки и сотни раз меньших, чем у соответствующих полигональных моделей. Однако использование SDF в интерактивных приложениях ограничено высокой вычислительной стоимостью классических методов рендеринга и отсутствием аппаратной поддержки в стандартном графическом конвейере. Попытки предварительной триангуляции перед визуализацией таких полей лишают преимуществ компактности хранения и гибкости изменения формы.

\label{ssec:actual}
Актуальность данной работы обусловлена необходимостью преодоления разрыва между эффективностью хранения неявных представлений и производительностью их визуализации. Появление технологии Mesh Shaders \cite{nvidia_mesh_shaders} открывает принципиально новые возможности для процедурной генерации геометрии, позволяя синтезировать полигональные фрагменты из функционального описания поверхности без записи в промежуточные буферы. 

Разработка метода, сочетающего компактность SDF с эффективностью аппаратного графического конвейера, позволит выйти на качественно новый уровень детализации в графических системах реального времени.

